% GENERATED FILE. Edit canonical lessons, then run npm run generate:books.
% canonical-source-hash: fnv1a64:fc4ee67c73a0ec22
% canonical-lessons: BN-C88-hajar, BN-C88-prithibi, BN-C88-rong, BN-C88-gan, BN-C88-chhuti

\chapter{Thousand, World, Colour, Song, Holiday}
\label{ch:bn-thousand-world-colour-song-holiday}
\hlchaptermodality{88}

\begin{chapteropening}
\textbf{By the end of this chapter:} \emph{I can say a thousand, the world, a colour, a song and a holiday.}

Complete the last lesson of chapter 88: I can say a thousand, the world, a colour, a song and a holiday.
\end{chapteropening}

\section[\texorpdfstring{\bn{হাজার} (hājār)}{হাজার (hājār)}]{\bn{হাজার} (hājār) — a thousand}

\label{lesson:BN-C88-hajar}



\begin{quote}
Before the new one: say the Bengali for fifty, then the Bengali for a hundred.
\end{quote}

\subsection*{You'll want to know: \bn{হাজার}}
\textbf{\bn{হাজার}} — \emph{hājār} — \textquotedblleft{}a thousand\textquotedblright{}.

\begin{cousinweb}
\textbf{\bn{হাজার}} (\emph{hājār}) is a loan, usually traced to Persian \emph{hazār}, \textquotedblleft{}thousand\textquotedblright{}.
\end{cousinweb}

\begin{grammarlens}[title={What you've built}]
One more word for this chapter.
\end{grammarlens}

\subsection*{Your turn}
\begin{itemize}
\raggedright
  \item \emph{Say it:} \emph{hājār}
  \item \emph{Say it:} \emph{hājār}, once more
  \item \emph{Say it:} say \emph{hājār}
  \item \emph{Recall:} say the Bengali for fifty, then the Bengali for a hundred, then say \emph{hājār} again
\end{itemize}

\begin{tcolorbox}[breakable,colback=teal!4,colframe=teal!35!black,title={Before you move on}]
What does \bn{হাজার} mean? (\textquotedblleft{}A thousand\textquotedblright{}.) Say it once more.
\end{tcolorbox}
\section[\texorpdfstring{\bn{পৃথিবী} (prithibi)}{পৃথিবী (prithibi)}]{\bn{পৃথিবী} (prithibi) — the world}

\label{lesson:BN-C88-prithibi}



\begin{quote}
Before the new one: say the Bengali for a hundred, then the Bengali for a thousand.
\end{quote}

\subsection*{You'll want to know: \bn{পৃথিবী}}
\textbf{\bn{পৃথিবী}} — \emph{prithibi} — \textquotedblleft{}the world\textquotedblright{}.

\begin{cousinweb}
\textbf{\bn{পৃথিবী}} (\emph{prithibi}) is Sanskrit \emph{pṛthivī}, \textquotedblleft{}the earth, the broad one\textquotedblright{}, taken over whole as a learned loan, a tatsama.
\end{cousinweb}

\begin{grammarlens}[title={What you've built}]
One more word for this chapter.
\end{grammarlens}

\subsection*{Your turn}
\begin{itemize}
\raggedright
  \item \emph{Say it:} \emph{prithibi}
  \item \emph{Say it:} \emph{prithibi}, once more
  \item \emph{Say it:} say \emph{prithibi}
  \item \emph{Recall:} say the Bengali for a hundred, then the Bengali for a thousand, then say \emph{prithibi} again
\end{itemize}

\begin{tcolorbox}[breakable,colback=teal!4,colframe=teal!35!black,title={Before you move on}]
What does \bn{পৃথিবী} mean? (\textquotedblleft{}The world\textquotedblright{}.) Say it once more.
\end{tcolorbox}
\section[\texorpdfstring{\bn{রং} (rông)}{রং (rông)}]{\bn{রং} (rông) — a colour}

\label{lesson:BN-C88-rong}



\begin{quote}
Before the new one: say the Bengali for a thousand, then the Bengali for the world.
\end{quote}

\subsection*{You'll want to know: \bn{রং}}
\textbf{\bn{রং}} — \emph{rông} — \textquotedblleft{}a colour\textquotedblright{}.

\begin{cousinweb}
\textbf{\bn{রং}} (\emph{rông}) matches both Sanskrit \emph{raṅga}, \textquotedblleft{}dye, colour\textquotedblright{}, under which CDIAL lists Bengali \emph{rāṅā}, \textquotedblleft{}red\textquotedblright{} (CDIAL 10560), and Persian \emph{rang}, \textquotedblleft{}colour\textquotedblright{}; the two may have run together.
\end{cousinweb}

\begin{grammarlens}[title={What you've built}]
One more word for this chapter.
\end{grammarlens}

\subsection*{Your turn}
\begin{itemize}
\raggedright
  \item \emph{Say it:} \emph{rông}
  \item \emph{Say it:} \emph{rông}, once more
  \item \emph{Say it:} say \emph{rông}
  \item \emph{Recall:} say the Bengali for a thousand, then the Bengali for the world, then say \emph{rông} again
\end{itemize}

\begin{tcolorbox}[breakable,colback=teal!4,colframe=teal!35!black,title={Before you move on}]
What does \bn{রং} mean? (\textquotedblleft{}A colour\textquotedblright{}.) Say it once more.
\end{tcolorbox}
\section[\texorpdfstring{\bn{গান} (gān)}{গান (gān)}]{\bn{গান} (gān) — a song}

\label{lesson:BN-C88-gan}



\begin{quote}
Before the new one: say the Bengali for the world, then the Bengali for a colour.
\end{quote}

\subsection*{You'll want to know: \bn{গান}}
\textbf{\bn{গান}} — \emph{gān} — \textquotedblleft{}a song\textquotedblright{}.

\begin{cousinweb}
\textbf{\bn{গান}} (\emph{gān}) is inherited from Sanskrit \emph{gāna}, \textquotedblleft{}singing\textquotedblright{}, and has come down with almost no change (CDIAL 4130).
\end{cousinweb}

\begin{grammarlens}[title={What you've built}]
One more word for this chapter.
\end{grammarlens}

\subsection*{Your turn}
\begin{itemize}
\raggedright
  \item \emph{Say it:} \emph{gān}
  \item \emph{Say it:} \emph{gān}, once more
  \item \emph{Say it:} say \emph{gān}
  \item \emph{Recall:} say the Bengali for the world, then the Bengali for a colour, then say \emph{gān} again
\end{itemize}

\begin{tcolorbox}[breakable,colback=teal!4,colframe=teal!35!black,title={Before you move on}]
What does \bn{গান} mean? (\textquotedblleft{}A song\textquotedblright{}.) Say it once more.
\end{tcolorbox}
\section[\texorpdfstring{\bn{ছুটি} (chhuṭi)}{ছুটি (chhuṭi)}]{\bn{ছুটি} (chhuṭi) — a holiday}

\label{lesson:BN-C88-chhuti}



\begin{quote}
Before the new one: say the Bengali for a colour, then the Bengali for a song.
\end{quote}

\subsection*{You'll want to know: \bn{ছুটি}}
\textbf{\bn{ছুটি}} — \emph{chhuṭi} — \textquotedblleft{}a holiday\textquotedblright{}.

\begin{cousinweb}
\textbf{\bn{ছুটি}} (\emph{chhuṭi}) is a tadbhava, a word inherited and reshaped by sound change. It comes from \emph{*kṣuṭyate}, \textquotedblleft{}is released\textquotedblright{}, a form CDIAL reconstructs from its descendants (CDIAL 3707).
\end{cousinweb}

\begin{grammarlens}[title={What you've built}]
One more word for this chapter.
\end{grammarlens}

\subsection*{Your turn}
\begin{itemize}
\raggedright
  \item \emph{Say it:} \emph{chhuṭi}
  \item \emph{Say it:} \emph{chhuṭi}, once more
  \item \emph{Say it:} say \emph{chhuṭi}
  \item \emph{Recall:} say the Bengali for a colour, then the Bengali for a song, then say \emph{chhuṭi} again
\end{itemize}

\begin{tcolorbox}[breakable,colback=teal!4,colframe=teal!35!black,title={Before you move on}]
What does \bn{ছুটি} mean? (\textquotedblleft{}A holiday\textquotedblright{}.) Say it once more.
\end{tcolorbox}
